\begin{abstract}
\noindent
The fiscal-multiplier literature treats government spending as an aggregate
resource flow, but the distributional incidence of public services varies
sharply across the income distribution: poor households complement public
provision, while wealthy households privately substitute. We formalize this
asymmetry by embedding \EC{} between public and private consumption into the
tractable heterogeneous-agent New Keynesian framework of \bilbiie{}. The
resulting model delivers \emph{two} sufficient statistics for fiscal
transmission -- the standard income-cyclicality $\chi$ and a new direct fiscal
channel $\xi$ -- that can be independently identified from household
consumption responses and reduced-form VAR evidence. Heterogeneous
internalization of public goods ($\theta^{H} \neq \theta^{S}$) induces
curvature heterogeneity that \emph{decouples the aggregate MPC from the
MPE}, providing a new route to the \citet{Auclert2024} trilemma through the
government-spending margin rather than labor-supply frictions. A
\emph{distributional dominance condition} links inequality, complementarity,
and the multiplier: fiscal policy is most effective precisely in economies
where it is most needed. Taking the model to a full HANK calibration
disciplined by US micro data, we find that plausibly heterogeneous
internalization raises impact multipliers to $1.70$--$1.84$, a $48$--$62\%$
amplification over the separable benchmark. Policy experiments reveal that
(i) targeting spending to complementary public services raises the multiplier
another $24\%$; (ii) austerity through public-service cuts is $44\%$ more
contractionary than conventional models predict and sharply regressive in
welfare; and (iii) in developing economies with high \HtM{} shares and
strongly regressive public-good quality, the impact multiplier reaches
$2.64$. The framework reframes fiscal policy as a question of
distributional design rather than aggregate size.
\end{abstract}
